\documentclass[11pt]{article}

% =====================================================================
%  The Spectral Geometry of Misalignment
%  Formal theory document for the alignment-geometry project.
%  Standard packages only; compiles with tectonic or latexmk -pdf.
% =====================================================================

\usepackage[margin=1in]{geometry}
\usepackage{amsmath,amssymb,amsthm,mathtools}
\usepackage{enumitem}
\usepackage[hidelinks]{hyperref}

\newtheorem{theorem}{Theorem}[section]
\newtheorem{proposition}[theorem]{Proposition}
\newtheorem{lemma}[theorem]{Lemma}
\newtheorem{corollary}[theorem]{Corollary}
\theoremstyle{definition}
\newtheorem{definition}[theorem]{Definition}
\newtheorem{assumption}{Assumption}
\theoremstyle{remark}
\newtheorem{remark}[theorem]{Remark}

\newcommand{\R}{\mathbb{R}}
\newcommand{\E}{\mathbb{E}}
\newcommand{\norm}[1]{\lVert #1 \rVert}
\newcommand{\inner}[2]{\langle #1, #2 \rangle}
\newcommand{\as}{\text{a.s.}}
\DeclareMathOperator{\Tr}{Tr}
\DeclareMathOperator{\rank}{rank}
\DeclareMathOperator{\diag}{diag}
\DeclareMathOperator*{\argmax}{arg\,max}
\newcommand{\dW}{\Delta W}
\newcommand{\Wb}{W_{\mathrm{b}}}
\newcommand{\Wm}{W_{\mathrm{m}}}
\newcommand{\MP}{\mathrm{MP}}
\newcommand{\TW}{\mathrm{TW}}

\title{\bfseries The Spectral Geometry of Misalignment}
\author{Aryan Gupta\\ \texttt{aryan.cs.app@gmail.com}}
\date{June 2026}

\begin{document}
\maketitle

\begin{abstract}
\noindent
Given a fine-tuned model and a matched benign control, we ask whether the
\emph{spectrum} of the weight update reveals a concentrated component associated
with a misaligned objective, without fitting the statistic on behavioral labels.
We model the weight increment $\dW = W_{\mathrm{ft}} - W_{\mathrm{base}}$ of a single
layer as an iid Gaussian rectangular matrix plus a deterministic fixed-rank signal,
and apply the additive rectangular singular-value transition of
Benaych-Georges and Nadakuditi \cite{benaych2012}. The classical
Baik--Ben~Arous--P\'ech\'e (BBP) spiked-covariance results
\cite{bbp2005,baik2006,paul2007} provide historical context but concern a
different, multiplicative Gaussian covariance model.
In the additive model, a direction of strength $\theta$ detaches from the
Marchenko--Pastur bulk \cite{marchenko1967} if and only if
$\theta>\sqrt\gamma$, where $\gamma$ is the layer aspect ratio. At fixed signal
energy $\tau$ and near-uniform allocation, a rank-$r$ update is detectable when
$r<r_\star=\tau/\sqrt\gamma$. Above threshold, the leading sample eigenvector
has computable overlap with the planted direction and provides a label-free
direction estimate. A known-scale Tracy--Widom null
\cite{tracywidom1996,johnstone2001} controls false positives under the Gaussian
asymptotic model; a fitted scale requires separate calibration. These conclusions
assume that a misaligned update concentrates its energy in fewer directions than
an energy-matched benign update. Whether real fine-tunes satisfy that assumption
is the principal empirical question. We give a label-free estimator and
characterize its main confounds and the limits of applying it to transformer
activations.
\end{abstract}

\section{Introduction}
\label{sec:intro}

Supervised white-box probes can read deception from residual-stream activations with
high accuracy when labeled examples from the relevant distribution are available
\cite{goldowskydill2025}. We consider the setting in which those labels and that
distribution are unavailable. In one four-team audit, the three teams with
model-and-data access succeeded and the API-only team failed; differing access
conditions and unusually searchable synthetic documents prevent attributing the
result to any single resource \cite{marks2025}.

We therefore analyze the weight increment
$\dW=W_{\mathrm{ft}}-W_{\mathrm{base}}$ directly. Its spectrum supplies a
same-recipe screening statistic that can be calibrated against a matched benign
control before behavioral validation.

Heavy-tailed self-regularization reads a label-free, data-free fingerprint of model
quality directly from weight spectra, including a ``bulk plus spikes'' regime
\cite{martin2021}. Spectral outlier analysis detects data poisoning from the singular
structure of learned representations \cite{tran2018}. Random-matrix methods have
also been applied to transformer internals for hallucination detection
\cite{ettori2026}, while related work connects weight spectra, activation
covariance, and fine-tuning in language models \cite{staats2024}. We ask whether
update concentration separates matched benign and misaligned weight updates.

\paragraph{Results and scope.}
The main results are:
\begin{enumerate}[leftmargin=1.5em,itemsep=2pt]
\item A random-matrix model of fine-tuning as a deterministic fixed-rank
perturbation of an iid Gaussian layer increment, with the diffuse part playing
the role of the Marchenko--Pastur bulk \cite{marchenko1967} and the additive
rectangular transition governed by \cite{benaych2012}
(Section~\ref{sec:model}).
\item A detectability theorem: a signal direction of strength $\theta$ is spectrally
visible exactly when $\theta > \sqrt{\gamma}$, with the eigenvalue location given in
closed form (Proposition~\ref{prop:detect}).
\item A \emph{rank-at-fixed-energy} discriminator: at matched
signal energy and equal or diffuse energy allocation, an update is detectable
exactly when its rank is below $r_\star = \tau/\sqrt{\gamma}$, so two
energy-matched fine-tunes of concentrated and diffuse rank are spectrally separated
(Proposition~\ref{prop:rank}).
\item A label-free identification result: above threshold the leading sample eigenvector
estimates the planted direction with computable overlap
(Proposition~\ref{prop:overlap}), and a known-scale Gaussian Tracy--Widom test
with ideal-null false-positive control (Proposition~\ref{prop:test}).
\item A subspace observable, the principal-angle distance between the leading eigenspaces
of matched models, intended to capture signal carried by eigenvector rotation rather than
eigenvalue magnitude (Definition~\ref{def:subspace}, Remark~\ref{rem:subspace}), with
surrogate calibration as a complement to scalar spectral summaries.
\item A characterization of the main failure modes (Section~\ref{sec:limits}),
including assumption violations for weights and activations, the limitations of
activation covariance, the relation between weight-update rank and activation
effective rank, and a matched-rank test of the hypothesis.
\end{enumerate}

\paragraph{Theoretical and empirical scope.}
Section~\ref{sec:results} establishes an ideal-model implication: if a misaligned
fine-tune deposits a rank-$r_{\mathrm m}$ signal and an equal-energy benign fine-tune
spreads its perturbation over a fixed rank $r_{\mathrm b}\gg r_{\mathrm m}$, then the
spectral test separates them at an explicit critical rank. The concentration hypothesis
is motivated by evidence that emergent misalignment is mediated by a single,
convergent linear direction
\cite{soligo2025}, a single rank-one adapter suffices to induce it \cite{turner2025}, and
emergent-misalignment weight deltas occupy a shared low-dimensional subspace across tasks
\cite{shared2025}. The remaining empirical quantity is the benign-update rank at
matched energy, measured as described in Section~\ref{sec:program}.

\paragraph{Spectral scope.}
The project name reflects an analogy with Fourier inspection of vision models. Weight
matrices lack a canonical periodic axis, so the analysis here uses singular values and
random-matrix theory. Fourier methods require additional group structure and fall outside
the present static-detection setting.

\section{Setup and notation}
\label{sec:setup}

Fix a layer with weight matrix $W \in \R^{p \times q}$, and write $W_{\mathrm{base}}$ for
its value in a reference model and $W_{\mathrm{ft}}$ for its value after fine-tuning. The
object of study is the increment
\[
  \dW \;=\; W_{\mathrm{ft}} - W_{\mathrm{base}} \;\in\; \R^{p \times q}.
\]
Without loss of generality take $q \le p$ (otherwise transpose), and set the aspect ratio
\[
  \gamma \;=\; q/p \;\in\; (0,1].
\]
We treat the $p$ rows of $\dW$ as $p$ observations in $\R^q$ and form the symmetric
positive semidefinite matrix
\[
  C \;=\; \frac{1}{p}\, \dW^{\!\top} \dW \;\in\; \R^{q \times q}.
\]
The nonzero eigenvalues of $C$ are the squared singular values of $\dW$ divided by
$p$. We write $\lambda_1(C) \ge \cdots \ge \lambda_q(C) \ge 0$ for its
eigenvalues and $\hat v_1, \dots, \hat v_q$ for the corresponding unit eigenvectors.

\paragraph{Orientation.}
The covariance $C=\dW^{\!\top}\dW/p$ is input-side: its eigenvectors are the right
singular vectors of $\dW$. For residual-stream interventions in the empirical paper we
use the output-side covariance $\widetilde C=\dW\dW^{\!\top}/q$, whose eigenvectors are the
left singular vectors. Transposition preserves singular values and algebraic rank, but
the displayed covariance normalizations differ: the nonzero eigenvalues of
$\widetilde C$ are $p/q$ times those of $C$. Its null edge and fluctuation scale must
therefore be rescaled accordingly. The dimensionless additive detectability condition is
equivalent after transposition. We state all scalar formulas in the $q\le p$, input-side
normalization, while behavioral interventions use output-side singular vectors.

\paragraph{Energy.}
The \emph{energy} of the update is its squared Frobenius norm,
$\norm{\dW}_F^2 = \Tr(\dW^{\!\top}\dW) = p\,\Tr(C)$. Energy is a weight-space quantity, and
we hold it fixed so that a benign and a misaligned update differ only in how their
movement is distributed across directions.

\paragraph{Matched models.}
The comparison uses three models sharing a base:
$W_{\mathrm{base}}$; a benign control $\Wb$ obtained by fine-tuning on an
alignment-neutral objective; and a misaligned model $\Wm$ obtained by an otherwise
identical recipe that installs the misaligned objective. Write
$\dW_{\mathrm b} = \Wb - W_{\mathrm{base}}$ and $\dW_{\mathrm m} = \Wm - W_{\mathrm{base}}$,
with associated $C_{\mathrm b}, C_{\mathrm m}$. The control is energy-matched,
$\norm{\dW_{\mathrm b}}_F = \norm{\dW_{\mathrm m}}_F$, which isolates the structure of
the update from its magnitude. Because any fine-tune can move the spectrum, comparison
with the base alone cannot attribute that movement to misalignment.

\section{Random-matrix preliminaries}
\label{sec:rmt}

The following standard results use the normalization of
Section~\ref{sec:setup}. Throughout, the asymptotic regime is
$p, q \to \infty$ with $q/p \to \gamma \in (0,1]$.

\subsection{The Gaussian bulk}

\begin{theorem}[Marchenko--Pastur law and Gaussian edge
\cite{marchenko1967,bai2010}]
\label{thm:mp}
Let $X \in \R^{p \times q}$ have iid $\mathcal N(0,\sigma^2)$ entries and let
$S = \tfrac1p X^{\!\top} X$. As $p,q \to \infty$ with $q/p \to \gamma
\in (0,1]$, the empirical spectral distribution of $S$ converges almost surely
to the Marchenko--Pastur law supported on $[\lambda_-,\lambda_+]$, where
\[
  \lambda_\pm=\sigma^2(1\pm\sqrt\gamma)^2.
\]
The largest eigenvalue also converges almost surely to
$\lambda_+=\sigma^2(1+\sqrt\gamma)^2$.
\end{theorem}

The upper edge is the asymptotic reference level for the ideal Gaussian null.
For finite matrices, an edge exceedance is descriptive until it is compared with
the null fluctuation scale of Theorem~\ref{thm:tw}.

\subsection{The additive rectangular spike model}

\begin{theorem}[Rectangular additive spike \cite{benaych2012}]
\label{thm:additive}
Let $N\in\R^{p\times q}$ have iid $\mathcal N(0,\sigma^2)$ entries, with
$q/p\to\gamma\in(0,1]$. Let $a\in\R^p$ and $b\in\R^q$ be deterministic unit
vectors, set $S=\beta ab^{\!\top}$, and define
$\theta=\beta^2/(p\sigma^2)$. For
$C=p^{-1}(N+S)^{\!\top}(N+S)$,
\[
  \lambda_1(C)\xrightarrow{\as}
  \begin{cases}
    \sigma^2(1+\theta)\left(1+\dfrac{\gamma}{\theta}\right),
      & \theta>\sqrt\gamma,\\[1.4ex]
    \sigma^2(1+\sqrt\gamma)^2, & \theta\le\sqrt\gamma.
  \end{cases}
\]
If $\hat b_1$ is the leading right singular vector of $N+S$, then
\[
  |\inner{\hat b_1}{b}|^2\xrightarrow{\as}
  \begin{cases}
    \dfrac{1-\gamma/\theta^2}{1+\gamma/\theta},
      & \theta>\sqrt\gamma,\\[1.4ex]
    0, & \theta\le\sqrt\gamma.
  \end{cases}
\]
\end{theorem}

The Gaussian specialization permits deterministic signal directions because the noise
is bi-orthogonally invariant. The general theory also treats fixed-rank perturbations
under explicit isotropy and edge assumptions. Growing-rank signals and correlated noise
require additional theory.

\subsection{Historical covariance analogue}

The classical Baik--Ben~Arous--P\'ech\'e literature
\cite{bbp2005,baik2006,paul2007} studies a different model: a low-rank
perturbation of a population covariance \cite{johnstone2001}. The
following Gaussian simple-spike statements explain the historical terminology and
provide the multiplicative covariance analogue.

\begin{theorem}[Classical spiked-covariance eigenvalue limit
\cite{baik2006}]
\label{thm:bbp}
Let the rows of $Y\in\R^{p\times q}$ be iid
$\mathcal N(0,\Sigma)$, where
$\Sigma=\sigma^2(I_q+\theta vv^{\!\top})$, $v$ is a deterministic unit vector,
and the population spike is simple. For $S_{\rm cov}=p^{-1}Y^{\!\top}Y$ and
$q/p\to\gamma\in(0,1]$,
\[
  \lambda_1(S_{\rm cov})\xrightarrow{\as}
  \begin{cases}
    \sigma^2(1+\theta)\left(1+\dfrac{\gamma}{\theta}\right),
      & \theta>\sqrt\gamma,\\[1.4ex]
    \sigma^2(1+\sqrt\gamma)^2, & \theta\le\sqrt\gamma.
  \end{cases}
\]
\end{theorem}

\begin{theorem}[Classical spiked-covariance eigenvector overlap
\cite{paul2007}]
\label{thm:overlap}
Under the Gaussian simple-spike conditions of Theorem~\ref{thm:bbp}, the
leading sample eigenvector $\hat v_1$ satisfies
\[
  |\inner{\hat v_1}{v}|^2\xrightarrow{\as}
  \begin{cases}
    \dfrac{1-\gamma/\theta^2}{1+\gamma/\theta},
      & \theta>\sqrt\gamma,\\[1.4ex]
    0, & \theta\le\sqrt\gamma.
  \end{cases}
\]
\end{theorem}

The additive and multiplicative models share these white-noise formulas while
retaining different probability laws.

\subsection{The known-scale null and a qualitative general bulk}

\begin{theorem}[Known-scale Gaussian Tracy--Widom edge
\cite{tracywidom1996,johnstone2001}]
\label{thm:tw}
For $S=p^{-1}N^{\!\top}N$ with iid real
$N_{ij}\sim\mathcal N(0,\sigma^2)$, known $\sigma^2$, and
$q/p\to\gamma\in(0,1]$,
\[
  \frac{\lambda_1(S)-\mu_{p,q}}{s_{p,q}}\Rightarrow\TW_1,
\]
where $\TW_1$ is the order-one Tracy--Widom law and
\[
  \mu_{p,q}=\frac{\sigma^2}{p}(\sqrt{p-1}+\sqrt q)^2,
  \qquad
  s_{p,q}=\frac{\sigma^2}{p}(\sqrt{p-1}+\sqrt q)
  \left(\frac1{\sqrt{p-1}}+\frac1{\sqrt q}\right)^{1/3}.
\]
The law was introduced for orthogonal ensembles by Tracy and Widom; Johnstone
established this real Wishart/PCA scaling.
\end{theorem}

\begin{remark}[Qualitative rectangular D-transform fallback]
\label{rem:rectangular-general}
For a non-white rectangular noise matrix with a compact limiting singular-value
law, no stray noise singular values, and the invariance or isotropic-direction
hypotheses of \cite{benaych2012}, a deterministic fixed-rank perturbation has a
qualitative outlier transition governed by the rectangular $D$-transform of the
noise singular-value law. The outlier location and threshold must be computed
from that law rather than from the Marchenko--Pastur closed forms. Applying this
result to checkpoint increments requires empirical checks of its invariance and
edge assumptions.
\end{remark}

\section{A spiked model of fine-tuning}
\label{sec:model}

We model the weight increment as a diffuse component plus a deterministic fixed-rank
component, with assumptions that give the diffuse term a Marchenko--Pastur bulk
\cite{marchenko1967}.

\begin{definition}[Diffuse-plus-signal decomposition]
\label{def:decomp}
Write the increment as
\[
  \dW \;=\; N \;+\; S, \qquad
  S \;=\; \sum_{i=1}^{r} \beta_i\, a_i b_i^{\!\top},
\]
where $N \in \R^{p\times q}$ is the diffuse component and $S$ is deterministic
with fixed rank $r=\rank(S)$ as $p,q\to\infty$. The vectors
$a_i\in\R^p$ and $b_i\in\R^q$ are orthonormal families and
$\beta_1\ge\cdots\ge\beta_r>0$. The dimensionless additive spike strengths are
$\theta_i=\beta_i^2/(p\sigma^2)$, and we consider sequences on which these
strengths have fixed positive limits. Here $\sigma^2$ is the per-entry variance
of $N$ from Assumption~\ref{ass:bulk}.
\end{definition}

\begin{assumption}[Ideal Gaussian diffuse bulk]
\label{ass:bulk}
The diffuse component has iid entries $N_{jk}\sim\mathcal N(0,\sigma^2)$.
\end{assumption}

Assumption~\ref{ass:bulk} supplies the closed-form thresholds. We apply it to the
increment because trained raw-weight spectra are often heavy-tailed or otherwise depart
from a pure Marchenko--Pastur law \cite{martin2021}. Gaussianity and isotropy of the
increment remain testable assumptions. Theorems~\ref{thm:mp}, \ref{thm:additive}, and
\ref{thm:tw} apply when those assumptions hold; Remark~\ref{rem:rectangular-general}
gives the qualitative non-white alternative.

\begin{assumption}[Structure hypothesis, to be tested]
\label{ass:structure}
A misalignment-relevant fine-tune deposits its signal in low rank, $r_{\mathrm m}$ small,
while a benign fine-tune of equal energy spreads a comparable Frobenius budget over higher
rank, $r_{\mathrm b} \gg r_{\mathrm m}$.
\end{assumption}

We treat Assumption~\ref{ass:structure} as an empirical hypothesis and condition the
separation results below on it.

\begin{remark}[Fixed-energy comparison]
\label{rem:energy}
By Definition~\ref{def:decomp}, the signal contributes
$\norm{S}_F^2 = \sum_i \beta_i^2 = p\,\sigma^2 \sum_i \theta_i$ to the energy. Holding the
signal energy fixed therefore fixes $\tau := \sum_i \theta_i$. The discriminator
of Section~\ref{sec:results} concerns how that $\tau$ is distributed across
$r$ directions, so the matched models of Section~\ref{sec:setup} are energy-matched.
\end{remark}

\section{Main results}
\label{sec:results}

\subsection{Detectability}

\begin{proposition}[Detectability and location of a single spike]
\label{prop:detect}
Under Assumption~\ref{ass:bulk} and the decomposition of Definition~\ref{def:decomp} with
a single signal direction of induced strength $\theta$, the leading eigenvalue of
$C = \tfrac1p \dW^{\!\top}\dW$ satisfies, as $p,q\to\infty$ with $q/p\to\gamma$,
\[
  \lambda_1(C) \;\xrightarrow{\as}\;
  \begin{cases}
    \sigma^2(1+\theta)\!\left(1+\dfrac{\gamma}{\theta}\right), & \theta > \sqrt{\gamma},\\[1.4ex]
    \sigma^2(1+\sqrt\gamma)^2, & \theta \le \sqrt{\gamma}.
  \end{cases}
\]
Consequently the signal is asymptotically visible in the spectrum, that is
$\lambda_1(C)$ exceeds the bulk edge by a deterministic gap, if and only if $\theta >
\sqrt{\gamma}$.
\end{proposition}

\begin{proof}
Apply Theorem~\ref{thm:additive} with
$\theta=\beta^2/(p\sigma^2)$ directly in the additive model.
\end{proof}

\subsection{The discriminator: rank at fixed energy}

For deterministic fixed $r$, the rectangular finite-rank theorem applies to each
planted singular direction. The detection of \emph{any} outlier is governed by
$\max_i\theta_i$, because the leading singular value tracks the strongest direction.

\begin{proposition}[Rank-at-fixed-energy discriminator]
\label{prop:rank}
For a fixed rank $r$ in the proportional limit, fix the signal energy through
$\tau = \sum_{i=1}^r \theta_i$. Among rank-$r$ signals of
this energy, the leading induced strength satisfies $\max_i \theta_i \ge \tau/r$, with
equality when the energy is spread equally, $\theta_i = \tau/r$. Therefore:
\begin{enumerate}[leftmargin=1.6em,itemsep=2pt]
\item An equal-energy, equal-strength rank-$r$ update is spectrally detectable
(Proposition~\ref{prop:detect}) if and only if
\[
  \frac{\tau}{r} > \sqrt{\gamma}
  \quad\Longleftrightarrow\quad
  r < r_\star, \qquad r_\star := \frac{\tau}{\sqrt{\gamma}}.
\]
\item Consequently, a concentrated update with $r_{\mathrm m} < r_\star$ produces a
supercritical spike above the bulk edge, while an energy-matched \emph{diffuse} update that
spreads its budget approximately uniformly over rank $r_{\mathrm b} \ge r_\star$ keeps every
direction subcritical and is asymptotically invisible. The test separates a concentrated
update from a diffuse one of equal energy; a high-rank update that nonetheless places a
supercritical share of its energy in a single direction is still detected.
\end{enumerate}
\end{proposition}

\begin{proof}
The inequality $\max_i \theta_i \ge \tau/r$ is the pigeonhole bound on a sum of $r$
nonnegative terms, with equality at the uniform allocation $\theta_i = \tau/r$. For the
uniform allocation, every spike has strength $\tau/r$, so by Proposition~\ref{prop:detect}
a detached eigenvalue exists if and only if $\tau/r > \sqrt\gamma$, equivalently $r <
\tau/\sqrt\gamma = r_\star$. For statement (2): when $r_{\mathrm m} < r_\star$ the leading
strength exceeds $\sqrt\gamma$ and a spike appears, while a uniform allocation over
$r_{\mathrm b} \ge r_\star$ puts every strength at or below $\sqrt\gamma$ and no eigenvalue
detaches. The pigeonhole bound is one-sided, so invisibility requires the diffuse,
near-uniform allocation rather than $r_{\mathrm b} \ge r_\star$ alone; a high-rank update
concentrated in one direction has $\max_i \theta_i > \sqrt\gamma$ and is detected.
\end{proof}

At fixed signal energy, rank controls detectability through the critical value
$r_\star=\tau/\sqrt\gamma$. Under Assumption~\ref{ass:structure}, the spectral test
separates matched updates whose concentrated and diffuse ranks lie on opposite sides of
$r_\star$. Section~\ref{sec:program} describes the corresponding measurement.

\begin{remark}[Uniform allocation minimizes detectability]
\label{rem:uniform}
Spreading energy uniformly minimizes the leading strength and is therefore the hardest
case to detect. If the uniform allocation is detectable, every allocation of the same
energy is detectable. Exact uniformity is unnecessary provided every benign spike remains
subcritical.
\end{remark}

\subsection{Critical rank for representative layers}
\label{sec:layers}

The critical rank $r_\star = \tau/\sqrt\gamma$ depends on the layer through its aspect
ratio $\gamma$ and on the fine-tune through the signal energy $\tau$. Table~\ref{tab:rstar}
reports $\gamma$ and
$r_\star$ for the released Llama-3-8B architecture (hidden width $4096$,
feed-forward width $14336$, 32 query heads, and eight key/value heads)
\cite{grattafiori2024}, at two illustrative signal energies.

\begin{table}[h]
\centering
\begin{tabular}{lccc}
\hline
Layer & shape $(p \times q)$ & $\gamma = q/p$ & $r_\star$ at $\tau = 4,\ 16$ \\
\hline
Attention $q/o$ projection & $4096 \times 4096$ & $1.00$ & $4.0,\ \ 16.0$ \\
Attention $k/v$ projection & $4096 \times 1024$ & $0.25$ & $8.0,\ \ 32.0$ \\
MLP up / gate projection & $14336 \times 4096$ & $0.286$ & $7.5,\ \ 29.9$ \\
MLP down projection & $14336 \times 4096$ & $0.286$ & $7.5,\ \ 29.9$ \\
\hline
\end{tabular}
\caption{Critical rank $r_\star = \tau/\sqrt\gamma$ across layer shapes. Shapes are written
as (rows $\times$ columns) after the transpose of Section~\ref{sec:setup} that places
$q \le p$, so $\gamma = q/p$ is the smaller dimension over the larger. An
equal-strength rank-$r$ signal is detectable in a layer exactly when
$r < r_\star$ for that layer.}
\label{tab:rstar}
\end{table}

The rectangular key/value and feed-forward maps have smaller $\gamma$ and therefore
higher critical ranks than the square query/output maps. The rank-one adapter that suffices
to induce emergent misalignment \cite{turner2025} sits at $r = 1$, far below $r_\star$ in
every layer, so the idealized equal-energy signal is well above threshold, whereas a
diffuse benign update would need to spread the same energy over at least $r_\star$
directions to stay hidden. Accordingly, the estimator in
Section~\ref{sec:estimator} scans layers and reports each standardized leading
eigenvalue; under the model, detection is easiest where $\gamma$ is small and
signal energy is large.

\subsection{Label-free identification}

Under the ideal model, a detected spike also yields an estimate of its planted direction.

\begin{proposition}[Label-free direction estimate]
\label{prop:overlap}
Under the conditions of Proposition~\ref{prop:detect} with $\theta > \sqrt\gamma$, the
leading eigenvector $\hat v_1$ of $C$ satisfies
\[
  |\inner{\hat v_1}{b}|^2 \;\xrightarrow{\as}\; \frac{1-\gamma/\theta^2}{1+\gamma/\theta} \;>\; 0,
\]
where $b$ is the planted signal direction. Thus $\hat v_1$ is asymptotically
nontrivially correlated with the perturbation direction, but is not a consistent
direction estimator at fixed $\theta$ and $\gamma$.
\end{proposition}

\begin{proof}
Immediate from the right-singular-vector statement of
Theorem~\ref{thm:additive}.
\end{proof}

Steering or ablating $\hat v_1$ tests its causal effect on behavior. This compares
the spectral estimate with the convergent activation direction reported for
emergent misalignment \cite{soligo2025}; the candidate itself is obtained without
alignment labels. Persona-vector results provide related evidence for
low-dimensional trait control \cite{chen2025}, but do not identify an
emergent-misalignment direction.

\subsection{A calibrated test}

\begin{proposition}[Known-scale level-$\alpha$ edge test]
\label{prop:test}
Let $H_0$ be the pure-noise null, $\dW = N$ with no signal (Assumption~\ref{ass:bulk} with
$\theta = 0$), and suppose $\sigma^2$ is known. Using the centering and scaling of
Theorem~\ref{thm:tw}, the test that rejects $H_0$ when
\[
  \frac{\lambda_1(C) - \mu_{p,q}}{s_{p,q}} \;>\; t_{1-\alpha},
  \qquad t_{1-\alpha} = \TW_1\text{ quantile},
\]
has asymptotic false-positive rate $\alpha$. Null eigenvalue fluctuations have
scale $p^{-2/3}$, while the critical window in spike strength around
$\theta=\sqrt\gamma$ has scale $p^{-1/3}$.
\end{proposition}

\begin{proof}
Under $H_0$ the leading eigenvalue obeys the Tracy--Widom limit of Theorem~\ref{thm:tw},
so the standardized statistic exceeds $t_{1-\alpha}$ with probability $\alpha$ in the
limit. For $\theta=\sqrt\gamma+\epsilon$, expansion of the outlier location in
Proposition~\ref{prop:detect} gives
\[
  \sigma^2(1+\theta)(1+\gamma/\theta)-\lambda_+
  =\frac{\sigma^2}{\sqrt\gamma}\epsilon^2+O(\epsilon^3).
\]
Matching this deterministic gap to the $p^{-2/3}$ edge scale gives
$\epsilon\asymp p^{-1/3}$. Calibrated power inside that critical window requires
a deformed-spike analysis beyond the pure-null Tracy--Widom approximation.
\end{proof}

Under the known-scale Gaussian null, Proposition~\ref{prop:test} gives asymptotic
type-I error control and the critical-window scale. Empirical use still requires
matched benign calibration and behavioral validation.

\subsection{A subspace observable for rotation}

Scalar spectral statistics depend only on the eigenvalues. Post-training can also rotate
representation geometry \cite{li2025geometry}, motivating a complementary eigenspace
observable.

\begin{definition}[Leading-subspace principal-angle distance
\cite{edelman1998}]
\label{def:subspace}
For two models with covariances $C, C'$, let $V_k, V_k'$ be the matrices of their top-$k$
eigenvectors and let $\rho_1 \ge \cdots \ge \rho_k$ be the singular values of $V_k^{\!\top}
V_k'$, equal to the cosines of the principal angles between the two subspaces. Define the
Grassmann distance
\[
  d_k(C, C') \;=\; \Big( \sum_{j=1}^k \arccos^2(\rho_j) \Big)^{1/2}.
\]
\end{definition}

\begin{remark}[Calibrating the rotation observable]
\label{rem:subspace}
The distance $d_k$ lacks Proposition~\ref{prop:test}'s scalar null. Under two independent
isotropic Gaussian null draws, the eigenvector bases are Haar-distributed
\cite{bai2010} and
independent leading directions are nearly orthogonal in high dimension. Matched
fine-tunes are correlated, so the independent-Haar null does not apply. We compare
the candidate with its energy-matched benign control and use the design-specific
resampling of Section~\ref{sec:permutation} only as a sensitivity analysis. The
resulting tail fraction is descriptive unless the transformations are
exchangeable under a stated null. An asymptotic law for $d_k$ would require a
joint model of matched increments.
\end{remark}

The subspace observable compares matched eigenspaces without behavioral labels and
measures reorientation of existing spectral structure.

\section{A label-free estimator}
\label{sec:estimator}

\subsection{The procedure}

\begin{enumerate}[leftmargin=1.6em,itemsep=3pt]
\item \textbf{Form the increment.} For each layer of interest compute $\dW = W_{\mathrm{ft}}
- W_{\mathrm{base}}$ and the covariance $C = \tfrac1p \dW^{\!\top}\dW$. Record the aspect
ratio $\gamma = q/p$.
\item \textbf{Estimate the bulk.} Fit the noise level $\sigma^2$ from the bulk of the
spectrum (for example by matching the median eigenvalue to the Marchenko--Pastur median),
rather than from the full trace, which a genuine outlier can inflate. Because
Proposition~\ref{prop:test} assumes a known scale, this plug-in estimate requires separate
calibration.
\item \textbf{Screen for an edge excursion.} Report the plug-in standardized leading
eigenvalue as a heuristic score, the ideal-model implied strength $\hat\theta$ from
Proposition~\ref{prop:detect}, and the descriptive number of MP-edge exceedances. The
exceedance count need not equal algebraic signal rank. A calibrated decision requires a
known scale or separate validation of the plug-in null.
\item \textbf{Match energy.} When comparing a misaligned candidate to a benign control,
rescale so $\norm{\dW_{\mathrm m}}_F = \norm{\dW_{\mathrm b}}_F$ before comparing spectra,
so the comparison is at fixed budget as Proposition~\ref{prop:rank} requires.
\item \textbf{Recover and test the direction.} Take $\hat v_1$ from
Proposition~\ref{prop:overlap}, map it back to a residual-stream direction, and test
causal coupling by steering or ablation. If the intervention changes behavior,
the spike is behaviorally relevant; other mechanisms may still contribute.
\item \textbf{Check rotation.} Compute the subspace distance of
Definition~\ref{def:subspace} between matched models to detect signal carried by
reorientation rather than magnitude.
\end{enumerate}

\subsection{Finite-size surrogate stress tests}
\label{sec:permutation}

Proposition~\ref{prop:test} provides asymptotic calibration for a Gaussian null with known
scale. Real increments are finite, non-Gaussian, correlated, and analyzed with a fitted
scale, so the plug-in Tracy--Widom statistic is a visibility heuristic. As a sensitivity
check, independently permuting all entries of $\dW$ preserves the global empirical entry
multiset while destroying row and column structure. Independent sign flips preserve
magnitudes and require a justified symmetric-sign null to preserve the signed marginal.
For $B$ surrogates, report the corrected tail fraction
\[
  \frac{1+\#\{b:\lambda_1(C^{(b)})\ge\lambda_1(C^{\rm obs})\}}{B+1}.
\]
Under a stated null that makes the transformations exchangeable, this quantity is a
randomization $p$-value; otherwise it is a descriptive surrogate tail fraction. Both
transformations discard optimizer, architectural, row, and column correlations.
The Llama-3-8B census therefore measures MP-visible structure but does not
establish significance under an empirical null.

\subsection{The activation-covariance instrument}
\label{sec:activation}

The same construction can be applied formally to the activation covariance
$\Sigma_{\mathrm{act}} = \tfrac1n \sum_{t} h_t h_t^{\!\top}$ of residual-stream states
$h_t$ over $n$ probe tokens for matched models. In one reported rank-one-LoRA
experiment, the layer-24 adapter $B$ vector and a mean-difference activation
direction have cosine near $0.04$, yet ablating that direction substantially
suppresses the measured behavior \cite{soligo2025}. This result motivates using
activation covariance for identification.

For detection, activation covariance falls outside
Propositions~\ref{prop:detect}--\ref{prop:test} because it is strongly anisotropic
before fine-tuning \cite{sun2024,gao2019,ethayarajh2019}. Assumption~\ref{ass:bulk}
and the resulting closed forms therefore do not apply. A generalized spiked
population model \cite{baiyao2012} offers an asymptotic analogue, but correlated
token states need not satisfy its independent-sample premise. Detection would
then depend on a less stable estimated deformed edge. Testing isotropy for weight
increments does not validate it for activation covariance.

\subsection{Weight-update rank and activation effective rank}
\label{sec:larf}

A recent result compares models fine-tuned on the top- versus bottom-ranked subsets
of nominally benign datasets and reports higher effective rank for the top-ranked,
safety-degrading subset when processing harmful prompts \cite{larf2025}. That result concerns
a different rank observable from the weight-increment rank studied here.

Effective rank \cite{roy2007} is $\exp(H)$ where $H = -\sum_i \pi_i \log \pi_i$ is the entropy of the
normalized spectrum $\pi_i = \lambda_i / \sum_j \lambda_j$. Adding a single dominant spike
to a covariance concentrates the distribution $\pi$, lowers $H$, and therefore
\emph{lowers} effective rank. Thus a low-rank misalignment signal lowers effective rank in
the relevant covariance.

The cited study measures inference-time activation effective rank after
fine-tuning on harmful \emph{content}, without energy control; our prediction
concerns the weight increment $\dW$. A concentrated update can still broaden the
range of behaviors expressed on triggering inputs, so the two observables need
not share a sign. Both should therefore be reported.

\subsection{Spike confounds}
\label{sec:confounds}

Several artifacts can produce a leading-eigenvalue excursion unrelated to misalignment.
\begin{itemize}[leftmargin=1.4em,itemsep=2pt]
\item \textbf{Outlier coordinates.} A few residual-stream coordinates with very large
variance \cite{sun2024} create a top eigenvalue unrelated to any fine-tune. Control by
per-coordinate standardization or robust covariance before spectral analysis, and by
working on $\dW$ where the effect is weaker.
\item \textbf{Energy leakage.} Without the energy matching of step~4, the statistic may
reflect training magnitude rather than misalignment.
\item \textbf{Aspect-ratio regime.} For very wide matrices or many probe tokens, $\gamma
\to 0$, the bulk collapses to a point and \emph{every} spike is nominally supercritical
($\sqrt\gamma \to 0$). Detectability is then trivial, and benign-versus-misaligned
separation depends entirely on the rank argument of Proposition~\ref{prop:rank}. Report
$\gamma$ explicitly and interpret the threshold only where it is nontrivial.
\item \textbf{Heavy-tailed bulk.} Many trained raw-weight spectra depart from a
pure Marchenko--Pastur law \cite{martin2021}. The Gaussian model is therefore
posed on $\dW$, but its bulk fit must still be checked. If it deviates, discard
the MP closed form and use Remark~\ref{rem:rectangular-general} only as a
qualitative guide.
\end{itemize}

\section{Empirical test}
\label{sec:program}

\paragraph{Matched-rank test.} Measure the rank of $\dW$ at matched energy for a benign control
and a misaligned model sharing a base, and locate both relative to $r_\star =
\tau/\sqrt\gamma$. Assumption~\ref{ass:structure} predicts $r_{\mathrm m} < r_\star \le
r_{\mathrm b}$. If the benign control is also low-rank at matched energy,
Proposition~\ref{prop:rank} gives no separation for that model pair. This
comparison directly tests Assumption~\ref{ass:structure}.

\paragraph{Substrate.} The matched comparison requires \emph{full} fine-tuning. A
rank-constrained adapter fixes the rank of $\dW$ by construction, precluding a rank
difference, and its increment has no bulk for a spike to cross.
Emergent-misalignment organisms exist as full fine-tunes and as rank-one adapters
\cite{betley2025,turner2025}; the rank-one adapters are useful for the identification and
steering checks of Proposition~\ref{prop:overlap}, but the detection claim of
Propositions~\ref{prop:detect}--\ref{prop:rank} must be evaluated on full fine-tunes with a
matched benign control trained under an identical recipe.

\paragraph{Prediction.} Under the rank-structure antecedent of
Assumption~\ref{ass:structure}, the ideal scalar prediction is a supercritical leading
eigenvalue in $\dW$ at a layer where the benign control remains subcritical. If benign
fine-tuning is also anisotropic, evidence must instead come from contrastive singular-vector
comparison, behavioral ablation, and calibrated same-recipe held-out screens. If
the measurements do not separate the models, the corresponding part of
Assumption~\ref{ass:structure} has failed.

\section{Limitations}
\label{sec:limits}

\begin{enumerate}[leftmargin=1.6em,itemsep=3pt]
\item \textbf{Rank separation at matched energy.} If
misaligned and benign fine-tunes are equally concentrated at matched energy, rank does not
discriminate and the method has no signal. Assumption~\ref{ass:structure} makes this
condition directly measurable.
\item \textbf{Activation anisotropy.} Activation covariance violates the
isotropic-bulk model used for $\dW$ and would require a deformed-bulk analysis
beyond the scope of this paper (Section~\ref{sec:activation}).
\item \textbf{Heavy tails on raw weights.} We place the Marchenko--Pastur null on the
diffuse part of the increment. Departures from MP in trained raw-weight spectra
\cite{martin2021} motivate avoiding a direct threshold on $W$ but do not establish the null
for increments.
\item \textbf{Narrow versus broad misalignment.} Narrow fine-tuning can induce
broad behavioral misalignment \cite{betley2025}; whether broader behavioral
effects correspond to more concentrated weight updates remains unknown.
\item \textbf{Internalized and deceptive objectives.} Belief--statement probes catch
instrumental deception but largely miss internalized objectives, which present no such gap
\cite{fanatics2026}. We target the internalized case, where behavioral probes are weakest
and no belief-gap signal is available. Any spectral signal would be structural and present
in the weights regardless of whether the model ever ``intends'' to deceive.
\item \textbf{Adaptive evasion.} An adaptive model could spread the update over
higher rank or move it into a benign-looking subspace, as demonstrated for other
latent-space monitors \cite{bailey2024}. Robustness to such evasion remains
untested.
\item \textbf{Finite-size and regime gaps.} The limits are asymptotic, whereas
real layers have finite $p,q$, and treating the rows of a single increment as
observations requires justification. Tracy--Widom calibration further assumes
Gaussianity and known scale. Without finite-matrix calibration, the fitted-scale
screen remains heuristic.
\end{enumerate}

\section{Related work}
\label{sec:related}

\paragraph{Spectral fingerprints of models.} Heavy-tailed self-regularization reads model
quality from weight spectra without data or labels, and names a ``bulk plus spikes'' phase
\cite{martin2021}. We apply this vocabulary to misalignment in the weight increment, using
the rectangular additive threshold rather than a power-law tail index. Work connecting
weight spectra, activation covariance, and transformer fine-tuning finds refinement along
small-singular-value directions \cite{staats2024}; our matched-energy comparison asks
whether benign and misaligned increments differ. Spectral fingerprints also support model
lineage detection designed to be invariant to fine-tuning \cite{ghost2026}. Random-matrix
descriptions of network spectra appear in analyses of loss-surface Hessians
\cite{pennington2017} and local spectral statistics of trained weight matrices, which remain
random-matrix-like even when the global density does not \cite{thamm2022}.

\paragraph{Spectral and random-matrix detection.} Spectral outlier scores for
data poisoning \cite{tran2018} and random-matrix hallucination detectors
\cite{ettori2026} operate per input. The present statistic operates per model and
uses rank at fixed energy.

\paragraph{The structure of emergent misalignment.} Emergent misalignment, in which narrow
fine-tuning induces broad misalignment \cite{betley2025}, is mediated by a single
convergent linear direction in activations \cite{soligo2025}, can be induced in clean
model organisms by a rank-one adapter \cite{turner2025}, is affected by persona
features \cite{wang2025persona,chen2025}, and has reported shared low-dimensional
parameter subspaces across tasks \cite{shared2025}. Together, these results
motivate Assumption~\ref{ass:structure} without establishing it. Under that
assumption, the rectangular additive model predicts an outlier transition
controlled by rank at fixed energy.
The single-direction phenomenon is part of a broader pattern in which
safety-relevant behaviors can be low-dimensional, as for refusal
\cite{arditi2024}; fine-tuning objectives can be optimized in low-intrinsic-dimensional
parameter subspaces \cite{aghajanyan2021}, and explicitly low-rank adapters are effective
\cite{hu2022}.

\paragraph{Geometry of fine-tuning and safety.} A complementary line analyzes how
fine-tuning breaks safety through low-rank curvature subspaces and argues that benign tasks
can also degrade alignment \cite{springer2026}, making the benign-side measurement
essential. Effective-rank and spectral-slope measurements likewise show that post-training
changes representation geometry \cite{li2025geometry}. Our analysis concerns detection
between matched models.

\paragraph{Behavioural and label-based detection.} Supervised probes detect deception at
high accuracy when labels and the distribution are available
\cite{goldowskydill2025}, and interpretability-driven audits surface hidden
objectives given sufficient access \cite{marks2025}. Unsupervised methods that search for
latent knowledge exist \cite{burns2023}, but arbitrary features can satisfy their
consistency constraints \cite{farquhar2023}. We focus on label-free detection
under distributional uncertainty; when suitable labels and a known distribution
are available, supervised probes should perform better.

\section{Conclusion}
\label{sec:conclusion}

At controlled energy, rank separates concentrated and diffuse updates in the fixed-rank
additive model, with critical rank $r_\star = \tau/\sqrt\gamma$ and a known-scale Gaussian
test. If misaligned fine-tunes are low-rank and matched benign fine-tunes are diffuse, the
spectral statistic separates them and estimates a candidate planted perturbation direction
without behavioral labels. Applying this result to real fine-tunes requires measuring
whether both updates lie on the predicted sides of $r_\star$, especially the matched benign
rank.

\begin{thebibliography}{99}

\bibitem{marchenko1967}
V.~A.~Mar\v{c}enko and L.~A.~Pastur,
``Distribution of eigenvalues for some sets of random matrices,''
\emph{Math.\ USSR-Sbornik} \textbf{1}(4):457--483, 1967.

\bibitem{bai2010}
Z.~Bai and J.~W.~Silverstein,
\emph{Spectral Analysis of Large Dimensional Random Matrices}, 2nd ed.,
Springer, 2010.

\bibitem{edelman1998}
A.~Edelman, T.~A.~Arias, and S.~T.~Smith,
``The geometry of algorithms with orthogonality constraints,''
\emph{SIAM J. Matrix Anal. Appl.} \textbf{20}(2):303--353, 1998.
\href{https://doi.org/10.1137/S0895479895290954}{doi:10.1137/S0895479895290954}.

\bibitem{johnstone2001}
I.~M.~Johnstone,
``On the distribution of the largest eigenvalue in principal components analysis,''
\emph{Ann.\ Statist.} \textbf{29}(2):295--327, 2001.

\bibitem{bbp2005}
J.~Baik, G.~Ben~Arous, and S.~P\'ech\'e,
``Phase transition of the largest eigenvalue for nonnull complex sample covariance matrices,''
\emph{Ann.\ Probab.} \textbf{33}(5):1643--1697, 2005.

\bibitem{baik2006}
J.~Baik and J.~W.~Silverstein,
``Eigenvalues of large sample covariance matrices of spiked population models,''
\emph{J.\ Multivariate Anal.} \textbf{97}(6):1382--1408, 2006.

\bibitem{paul2007}
D.~Paul,
``Asymptotics of sample eigenstructure for a large dimensional spiked covariance model,''
\emph{Statist.\ Sinica} \textbf{17}:1617--1642, 2007.

\bibitem{benaych2012}
F.~Benaych-Georges and R.~R.~Nadakuditi,
``The singular values and vectors of low rank perturbations of large rectangular random matrices,''
\emph{J.\ Multivariate Anal.} \textbf{111}:120--135, 2012.

\bibitem{tracywidom1996}
C.~A.~Tracy and H.~Widom,
``On orthogonal and symplectic matrix ensembles,''
\emph{Comm.\ Math.\ Phys.} \textbf{177}(3):727--754, 1996.

\bibitem{baiyao2012}
Z.~Bai and J.~Yao,
``On sample eigenvalues in a generalized spiked population model,''
\emph{J.\ Multivariate Anal.} \textbf{106}:167--177, 2012.

\bibitem{pennington2017}
J.~Pennington and Y.~Bahri,
``Geometry of neural network loss surfaces via random matrix theory,''
\emph{ICML}, PMLR~70:2798--2806, 2017.

\bibitem{martin2021}
C.~H.~Martin and M.~W.~Mahoney,
``Implicit self-regularization in deep neural networks: evidence from random matrix theory
and implications for learning,''
\emph{J.\ Mach.\ Learn.\ Res.} \textbf{22}(165):1--73, 2021.

\bibitem{thamm2022}
M.~Thamm, M.~Staats, and B.~Rosenow,
``Random matrix analysis of deep neural network weight matrices,''
\emph{Phys.\ Rev.\ E} \textbf{106}:054124, 2022.

\bibitem{staats2024}
M.~Staats, M.~Thamm, and B.~Rosenow,
``Small singular values matter: a random matrix analysis of transformer models,''
\emph{NeurIPS}, vol.~38, 2025. arXiv:2410.17770.

\bibitem{roy2007}
O.~Roy and M.~Vetterli,
``The effective rank: a measure of effective dimensionality,''
\emph{European Signal Processing Conference (EUSIPCO)}, 606--610, 2007.

\bibitem{grattafiori2024}
A.~Grattafiori et al.,
``The Llama 3 herd of models,''
arXiv preprint arXiv:2407.21783, 2024.

\bibitem{betley2025}
J.~Betley, D.~C.~H.~Tan, N.~Warncke, A.~Sztyber-Betley, X.~Bao, M.~Soto, N.~Labenz, and O.~Evans,
``Emergent misalignment: narrow finetuning can produce broadly misaligned LLMs,''
\emph{ICML}, PMLR~267:4043--4068, 2025. arXiv:2502.17424.

\bibitem{turner2025}
E.~Turner, A.~Soligo, M.~Taylor, S.~Rajamanoharan, and N.~Nanda,
``Model organisms for emergent misalignment,''
arXiv:2506.11613, 2025.

\bibitem{soligo2025}
A.~Soligo, E.~Turner, S.~Rajamanoharan, and N.~Nanda,
``Convergent linear representations of emergent misalignment,''
arXiv:2506.11618, 2025.

\bibitem{wang2025persona}
M.~Wang, T.~Dupr\'e~la~Tour, O.~Watkins, A.~Makelov, R.~A.~Chi, S.~Miserendino,
J.~Wang, A.~Rajaram, J.~Heidecke, T.~Patwardhan, and D.~Mossing,
``Persona features control emergent misalignment,''
\emph{ICLR}, 2026. arXiv:2506.19823.

\bibitem{chen2025}
R.~Chen, A.~Arditi, H.~Sleight, O.~Evans, and J.~Lindsey,
``Persona vectors: monitoring and controlling character traits in language models,''
arXiv:2507.21509, 2025.

\bibitem{shared2025}
D.~A.~R.~Arturi, E.~Zhang, A.~Ansah, K.~Zhu, A.~Panda, and A.~Balwani,
``Shared parameter subspaces and cross-task linearity in emergently misaligned behavior,''
arXiv:2511.02022, 2025.

\bibitem{arditi2024}
A.~Arditi, O.~Obeso, A.~Syed, D.~Paleka, N.~Panickssery, W.~Gurnee, and N.~Nanda,
``Refusal in language models is mediated by a single direction,''
\emph{NeurIPS}, 2024. arXiv:2406.11717.

\bibitem{aghajanyan2021}
A.~Aghajanyan, S.~Gupta, and L.~Zettlemoyer,
``Intrinsic dimensionality explains the effectiveness of language model fine-tuning,''
\emph{ACL-IJCNLP}, 2021. arXiv:2012.13255.

\bibitem{hu2022}
E.~J.~Hu, Y.~Shen, P.~Wallis, Z.~Allen-Zhu, Y.~Li, S.~Wang, L.~Wang, and W.~Chen,
``LoRA: low-rank adaptation of large language models,''
\emph{ICLR}, 2022. arXiv:2106.09685.

\bibitem{springer2026}
M.~Springer, C.~P.~Lee, B.~Metevier, J.~Castleman, B.~Turbal, H.~Jung,
Z.~Shen, and A.~Korolova,
``The geometry of alignment collapse: when fine-tuning breaks safety,''
arXiv:2602.15799, 2026. Preprint.

\bibitem{larf2025}
H.~Li, L.~Li, Z.~Lu, X.~Wei, R.~Li, J.~Shao, and L.~Sha,
``Layer-aware representation filtering: purifying finetuning data to preserve LLM safety alignment,''
\emph{EMNLP}, 8030--8050, 2025. doi:10.18653/v1/2025.emnlp-main.406.

\bibitem{li2025geometry}
M.~Z.~Li, K.~K.~Agrawal, A.~Ghosh, K.~K.~Teru, A.~Santoro, G.~Lajoie, and B.~A.~Richards,
``Tracing the representation geometry of language models from pretraining to post-training,''
\emph{NeurIPS}, vol.~38, 2025. arXiv:2509.23024.

\bibitem{ettori2026}
D.~Ettori,
``Spectral Geometry for Deep Learning: Compression and Hallucination Detection via Random
Matrix Theory,''
arXiv:2601.17357, 2026. MS thesis, University of Illinois Chicago,
defended November 21, 2025.

\bibitem{tran2018}
B.~Tran, J.~Li, and A.~M\k{a}dry,
``Spectral signatures in backdoor attacks,''
\emph{NeurIPS}, 2018. arXiv:1811.00636.

\bibitem{goldowskydill2025}
N.~Goldowsky-Dill, B.~Chughtai, S.~Heimersheim, and M.~Hobbhahn,
``Detecting strategic deception using linear probes,''
arXiv:2502.03407, 2025.

\bibitem{marks2025}
S.~Marks, J.~Treutlein, T.~Bricken, J.~Lindsey, J.~Marcus,
S.~Mishra-Sharma, D.~Ziegler, E.~Ameisen, J.~Batson, T.~Belonax,
S.~R.~Bowman, S.~Carter, B.~Chen, H.~Cunningham, C.~Denison,
F.~Dietz, S.~Golechha, A.~Khan, J.~Kirchner, J.~Leike, A.~Meek,
K.~Nishimura-Gasparian, E.~Ong, C.~Olah, A.~Pearce, F.~Roger,
J.~Salle, A.~Shih, M.~Tong, D.~Thomas, K.~Rivoire, A.~Jermyn,
M.~MacDiarmid, T.~Henighan, and E.~Hubinger,
``Auditing language models for hidden objectives,''
arXiv:2503.10965, 2025.

\bibitem{burns2023}
C.~Burns, H.~Ye, D.~Klein, and J.~Steinhardt,
``Discovering latent knowledge in language models without supervision,''
\emph{ICLR}, 2023. arXiv:2212.03827.

\bibitem{farquhar2023}
S.~Farquhar, V.~Varma, Z.~Kenton, J.~Gasteiger, V.~Mikulik, and R.~Shah,
``Challenges with unsupervised LLM knowledge discovery,''
arXiv:2312.10029, 2023.

\bibitem{bailey2024}
L.~Bailey, A.~Serrano, A.~Sheshadri, M.~Seleznyov, J.~Taylor,
E.~Jenner, J.~Hilton, S.~Casper, C.~Guestrin, and S.~Emmons,
``Obfuscated activations bypass LLM latent-space defenses,''
\emph{ICLR}, 2026. arXiv:2412.09565.

\bibitem{fanatics2026}
K.~Haralambiev,
``Why safety probes catch liars but miss fanatics,''
arXiv:2603.25861, 2026. Preprint.

\bibitem{sun2024}
M.~Sun, X.~Chen, J.~Z.~Kolter, and Z.~Liu,
``Massive activations in large language models,''
\emph{COLM}, 2024. arXiv:2402.17762.

\bibitem{gao2019}
J.~Gao, D.~He, X.~Tan, T.~Qin, L.~Wang, and T.-Y.~Liu,
``Representation degeneration problem in training natural language generation models,''
\emph{ICLR}, 2019. arXiv:1907.12009.

\bibitem{ethayarajh2019}
K.~Ethayarajh,
``How contextual are contextualized word representations? Comparing the geometry of BERT,
ELMo, and GPT-2 embeddings,''
\emph{EMNLP-IJCNLP}, 2019.

\bibitem{ghost2026}
S.~Wang, Z.~Ma, X.~Li, and Z.~Li,
``Ghost in the Transformer: Detecting Model Reuse with Invariant Spectral Signatures,''
\emph{Proc. AAAI Conf. Artif. Intell.} \textbf{40}(40):33648--33656, 2026.
doi:10.1609/aaai.v40i40.40654. arXiv:2511.06390.

\end{thebibliography}

\end{document}
